% !TEX TS-program = pdflatex
% David - Harjoittelija CV (suomi)

\documentclass[a4paper,11pt]{article}

\IfFileExists{cv-style.sty}{\usepackage{cv-style}}{\usepackage{../../shared/cv-style}}

\cvlang{fi}
\cvrole{Harjoittelija}

\begin{document}
\small
\setlength{\cvlabelw}{3.9cm}
\titlespacing{\section}{0pt}{0.55em}{2pt}

\begin{cvheader}
\par{\centering
		{\Huge \textsc{David Adegbola}}\\[0.4em]
		{\Large Harjoittelija}
	\bigskip\par}

\hrule
\vspace{0.5em}
\cvline[2.4cm]{\cvt{location}:}{Joensuu, Suomi}
\cvline[2.4cm]{\cvt{phone}:}{+358 46 800 7210}
\cvline[2.4cm]{\cvt{email}:}{\href{mailto:adavidadegbola@gmail.com}{adavidadegbola@gmail.com}}
\end{cvheader}

\section{\cvt{summary}}
\begin{tabular}{p{0.9\textwidth}}
 Hei! Olen David, 24-vuotias päätoiminen tutkinto-opiskelija Suomessa: opiskelen kestävää metsäbiotaloutta Itä-Suomen yliopistossa. Taustani on tietotekniikassa, sillä opiskelin tietotekniikkaa VAMKissa ja suoritin OND-tutkinnon tietojenkäsittelystä. UPM Kaukaan sellutehtaalla työskentelin suurella teollisuustyömaalla. Työnantaja arvioi työni kiitettäväksi ja kuvasi minua halukkaaksi oppimaan ja ottamaan vastuuta.\\\\

 Vaasa Startup Factoryssa tein töitä tiimissä: haastattelimme käyttäjiä, soitimme yrityksille ja esittelimme tuotteemme yleisölle. Puhun sujuvaa englantia ja voin työskennellä Vaasassa. Suomen kielen taitoni on vielä perustasolla, ja opiskelen sitä aktiivisesti.
\end{tabular}

\section{\cvt{education}}
\cvline{2026--nykyhetki:}{B.Sc, kestävä metsäbiotalous, Itä-Suomen yliopisto (UEF), päätoiminen opiskelija}
\cvline{2024--2026:}{B.Eng-opinnot, tietotekniikka, Vaasan ammattikorkeakoulu (VAMK), kesken; opinnoissa Microsoft Active Directory}
\cvline{2019--2022:}{OND, tietojenkäsittely, The Polytechnic, Ibadan, Nigeria}

\section{\cvt{skills}}
\cvline{Tietotekniikka:}{Microsoft Office; Microsoft Active Directory (VAMKin opinnot)}
\cvline{Selvitystyö:}{Käyttäjähaastattelut ja yrityksille soittaminen}
\cvline{Työtavat:}{Tiimityö; esiintyminen; viestintä ja yhteistyö}
\cvline{Turvallisuus:}{Teollisuustyömaan turvallisuusohjeiden noudattaminen}
\cvline{Haluan oppia:}{Tehtävässä käytettävät data-analyysityökalut}
\cvline{\cvt{cards}:}{Työturvallisuuskortti}
\cvline{\cvt{languages}:}{Englanti (sujuva); suomi (perustaso, opiskelen aktiivisesti); yoruba (äidinkieli)}

\section{\cvt{experience}}
\begin{cventry}{6/2026--8/2026}{\job{Founder in Residence}{Vaasa Startup Factory}, Vaasa}
  Selvitin sähkön hinnan vaikutusta tuotantokustannuksiin käyttäjähaastatteluilla ja yrityksille soittamalla.\par
  Kehitimme tiimini kanssa hetkin.app-palvelun ja esittelin sen yleisölle.
\end{cventry}
\begin{cventry}{18.8.2025--7.10.2025}{\job{Asentaja}{Axioma Servisas UAB}}
  UPM Kaukaan sellutehdas, Lappeenranta: sähkösuodattimen uudelleenrakennushanke.\par
  Noudatin tehtaan tiukkoja turvallisuusohjeita, pidin työmaan siistinä ja seurasin työkaluja ja materiaaleja.\par
  Työnantajan arvio: kiitettävä. Vastuuntuntoinen ja tunnollinen työntekijä, halukas oppimaan ja ottamaan vastuuta.
  \sotagrow{teollisuustyömaa,työturvallisuus,vastuullisuus,tiimityö,oppimishalu}
\end{cventry}
\begin{cventry}{11/2024--8/2026}{\textbf{AV-vapaaehtoinen}, Vaasa City Church, Vaasa: hoidin ääni- ja videotekniikan joka viikonloppu.}
\end{cventry}
\begin{cventry}{4/2025--8/2025}{\textbf{Lehdenjakaja}, Finnbox ja SSM, Vaasa: jakelureitit yksin aikaisin aamulla.}
\end{cventry}
\begin{cventry}{9/2016--9/2020}{Myynti- ja asiakaspalvelutyötä Nigeriassa (Desire Feeds, Bunmi Pastries): kassan käsittely.}
\end{cventry}

\section{\cvt{hobbies}}
Kuntosali, asioiden korjaaminen ja valokuvaus.

\section{\cvt{references}}
\cvt{refsavailable}

\end{document}
